% Dataset section. Main text is deliberately short: design and composition
% here, procedure and validity evidence in Appendix A (sections/dataset_details).
%
% Dataset: dataset/multi_domain_dataset/. Every number below was read from the
% data files themselves, not from any summary document:
%   */dataset_manifest.json              counts, label distributions, date range
%   mixed/idea_novelty_dataset.jsonl     per-domain novelty distribution (the
%       dataset carries a native `domain` field; nothing is recovered by matching)
%   mixed/pair_staged/*.jsonl            pair split sizes
% Idea-level split sizes follow from applying the pair_staged date cuts
% (< 202502, 202502--202510, >= 202511) to target_date; they are reproducible
% from idea_novelty_dataset.jsonl alone.
%
% OPEN: Section 5 and Section 6 still carry numbers from the earlier 25,435-row
% build (3,816 test / 3,177 domain-identified / 17,804-3,815-3,816 splits) and
% need replacing with results measured on this dataset.

\section{Dataset}
\label{sec:dataset}

Novelty can only be assessed relative to earlier work. Each evaluation example
must therefore specify a target idea, the prior ideas available when it was
proposed, and the contribution that remains after accounting for those priors.
We construct \textsc{IdeaDelta} to provide these three elements.

\subsection{Task and labels}
\label{sec:dataset-task}

The unit of annotation is a research \emph{idea}: one innovation claim extracted
from a paper. For example, a paper that introduces both a new objective and a
new architecture contributes two ideas, which are evaluated separately. Let
$x$ denote a target idea published at time $t_x$. Its admissible priors are
$\mathcal{P}_x=\{p_k\mid t_{p_k}<t_x\}$. This temporal constraint is essential:
an idea published after the target cannot serve as evidence about the target's
novelty, regardless of its semantic similarity. We annotate each target--prior
pair with one of four coverage levels and each target with an ordered novelty
grade from \textsc{low} to \textsc{high}.

Rather than assign novelty as a single overall impression, we derive it from
four annotated quantities:
\begin{equation}
 \mathrm{raw}(x)=\bigl(1-C_x\bigr)\cdot\Delta_x\cdot E_x\cdot
                 \bigl(1-\lambda S_x\bigr),
 \qquad \lambda=0.45,
 \label{eq:novelty-rule}
\end{equation}
where $C_x$ measures how much of the target is jointly covered by the retained
priors, $\Delta_x$ measures the substantive contribution that remains,
$E_x$ measures whether the retrieved evidence is sufficient for assessment,
and $S_x$ measures how much of the claimed difference is merely surface-level.
We convert the resulting score into the four ordered novelty grades. The rule
assigns lower novelty when prior work already covers the target, when the
remaining difference is not substantive, when the available evidence is
insufficient, or when the difference consists mainly of renaming or other
surface changes.

\subsection{Construction}
\label{sec:dataset-construction}

\paragraph{Paper selection and idea extraction.}
We collect arXiv papers from graph machine learning, computer vision, and time
series. We prioritize papers that were subsequently published, especially at
major venues. An LLM extracts innovation claims from the title, abstract, and
introduction. A second pass removes statements that do not describe a
substantive research contribution, including open-source announcements, bare
state-of-the-art claims, and engineering deliverables.

\paragraph{Prior retrieval.}
For each target, nearest-neighbor search over idea embeddings retrieves
candidate priors published before the target. An LLM adjudicates candidates in
the uncertain similarity range. We also perform a cross-type retrieval pass so
that a target can be linked to relevant prior work of another contribution type.
All target--prior pairs remain within the same research domain.

\paragraph{Annotation.}
The same judge, \texttt{deepseek-v4-pro}, annotates all stages. It first assigns
pairwise coverage labels, with ambiguous cases judged from their technical
content. It then estimates joint coverage, the remaining substantive
contribution, evidence sufficiency, and surface-level difference for each
target. Equation~(\ref{eq:novelty-rule}) converts these quantities into the
idea-level novelty grade. The labels are model-generated, and we manually
inspect samples from each stage. Appendix~\ref{app:dataset} describes the full
pipeline.

\subsection{Composition}
\label{sec:dataset-composition}

\textsc{IdeaDelta} contains $21{,}111$ annotated ideas and $92{,}876$ annotated
target--prior pairs, with target dates ranging from 2019-01 to 2026-09
(Table~\ref{tab:dataset-composition}). Each idea records its source domain, and
the \emph{All} column is the union of the three domain subsets rather than an
additional collection. We order the train, validation, and test splits by date,
with cut points at 2025-02 and 2025-11. The resulting splits contain
$59{,}269/16{,}436/17{,}171$ pairs and $13{,}815/3{,}944/3{,}352$ ideas.

\begin{table}[t]
\centering
\small
\setlength{\tabcolsep}{6pt}
\caption{Dataset composition. Counts of ideas and of target--prior pairs, and the
idea-level novelty distribution. Pairs are formed within a source domain only.}
\label{tab:dataset-composition}
\begin{tabular}{@{}lrrrr@{}}
\toprule
& Graph ML & Computer vision & Time series & All \\
\midrule
Ideas & 5{,}225 & 8{,}291 & 7{,}595 & 21{,}111 \\
Pairs & 21{,}183 & 36{,}238 & 35{,}455 & 92{,}876 \\
\addlinespace
\textsc{low} & 819 & 901 & 1{,}724 & 3{,}444 \\
\textsc{weak} & 1{,}240 & 1{,}711 & 1{,}614 & 4{,}565 \\
\textsc{moderate} & 2{,}247 & 2{,}777 & 2{,}173 & 7{,}197 \\
\textsc{high} & 919 & 2{,}902 & 2{,}084 & 5{,}905 \\
\bottomrule
\end{tabular}
\end{table}

Because the labels are model-generated, we further test whether they behave as
novelty labels should. Appendix~\ref{app:dataset-validity} examines whether
metadata can predict the labels, whether the retrieved prior content matters,
and whether the assigned novelty changes appropriately when the evidence is
modified.
